\section{逆関数による置換}\label{節:逆向きの置換}

これまでは,積分される式の中に $g(x)$ とその微分 $g'(x)$ の組を探した.では,\,$\sqrt{x}$ のように根号があって,その微分が見当たらない場合はどうすればよいだろうか.この場合は,根号などのまとまりを $t=g(x)$ と置いてから,元の変数を $x=h(t)$ と解き直し,\,$x$ と $\dd x$ をすべて $t$ で書き直す.根号を含まない $t$ の式,たとえば前章の割り算や部分分数で積分できる形に帰着させることを目標とする.

ここで1つ断りをいれておこう.\,$t=g(x)$ から $x$ を $t$ の1つの関数 $x=h(t)$ として表せるのは,\,$g$ が狭義単調な区間に限る.本節では,そのような区間を選んで置換する.その理由は例題の後で述べる.

根拠は,\sectionref{節:合成形の置換}の公式を $t$ の関数 $h$ に使うことである.\,$F'=f$ ならば
\[
 \frac{\dd}{\dd t}F(h(t))=f(h(t))h'(t)
\]
だから,
\[
 \int f(h(t))h'(t)\dd t=F(h(t))+C
\]
となる.ここで $t=g(x)$ を戻すと,右辺は $F(h(g(x)))$ である.これが求める $F(x)$ になるためには,\,$h(g(x))=x$,すなわち $h$ が $g$ の逆関数であればよい.したがって,\,$x=h(t)$,\,$\dd x=h'(t)\dd t$ と書き直して $t$ で積分し,最後に $t=g(x)$ を戻せばよいのだ.

たとえば根号 $\sqrt{ax+b}$ $(a\ne0)$ を含む場合は,\,$t=\sqrt{ax+b}$ と置く.\,$t\ge0$ であり,
\[
 x=\frac{t^2-b}a,\qquad \dd x=\frac{2t}a\dd t
\]
となるから,根号が $t$ に変わる.

\begin{bookpoint}
\textbf{POINT}\par
根号などを $t$ と置いて $x=h(t)$ と解き直し,\,$\dd x=h'(t)\dd t$ として,被積分関数と $\dd x$ をすべて $t$ で書き直す.
\end{bookpoint}

\begin{example}{根号への置換}
$x>0$ において,\,$\int\frac{\sqrt{x}}{1+\sqrt{x}}\dd x$ を求めよ.
\end{example}

\paragraphbf{方針}
$t=\sqrt{x}$ と置いて $x=t^2$ と解き直し,\,$t$ の分数式にしてから割り算をする.

\paragraphbf{解答}
$x-2\sqrt{x}+2\log(1+\sqrt{x})+C$

\paragraphbf{解法}
$t=\sqrt{x}$ と置けば,\,$t>0$ で $x=t^2$,\,$\dd x=2t\dd t$ となる.分子だけでなく分母と $\dd x$ も書き換えて
\begin{align*}
 \int\frac{\sqrt{x}}{1+\sqrt{x}}\dd x
 &=\int\frac{t}{1+t}\cdot2t\dd t\\
 &=2\int\frac{t^2}{t+1}\dd t
\end{align*}
となる.根号はなくなったが分数が残った.ここで前章の割り算を使い,\,$t^2=(t+1)(t-1)+1$ と分ける.
\begin{align*}
 2\int\frac{t^2}{t+1}\dd t
 &=2\int\left(t-1+\frac1{t+1}\right)\dd t\\
 &=t^2-2t+2\log(t+1)+C
\end{align*}
$t>0$ では $t+1>0$ なので,対数の絶対値を外せる.\,$t=\sqrt{x}$ を戻すと,答えは
\[
 x-2\sqrt{x}+2\log(1+\sqrt{x})+C
\]
である.微分して確かめると
\begin{align*}
 1-\frac1{\sqrt{x}}+\frac1{(1+\sqrt{x})\sqrt{x}}
 &=1-\frac1{1+\sqrt{x}}\\
 &=\frac{\sqrt{x}}{1+\sqrt{x}}
\end{align*}
となる.置換が積分を1回で終わらせなくても,根号を有理式へ変えるという目的は達している.

\paragraphbf{枝と単調性}
この例で $x=t^2$ と書けたのは,\,$t=\sqrt{x}>0$ と決めていたからである.逆に $t=x^2$ から $x$ を求めるなら,\,$x=\sqrt{t}$ と $x=-\sqrt{t}$ の2つがある.\,$x>0$ の区間では前者,\,$x<0$ の区間では後者を使う.このように,逆関数を1つに決めるために選ぶ範囲を\textbf{枝}と呼ぶことがある.

区間上の連続関数 $g$ が逆関数を持つことと,\,$g$ が狭義単調であることは同値である.ここで狭義単調とは,\,$x_1<x_2$ ならつねに $g(x_1)<g(x_2)$ となるか,つねに $g(x_1)>g(x_2)$ となることをいう.したがって,\,$t=g(x)$ から1つの関数 $x=h(t)$ を使って戻したいなら,そのような区間を選ばなければならない.さらに逆関数の微分を $h'(t)=\frac1{g'(h(t))}$ と書く場合には,分母が $0$ でないことも必要である.

これに対し,\sectionref{節:合成形の置換}の公式では逆関数を作っていない.\,$F(g(x))$ をそのまま微分したので,単調性を仮定しなかった.定積分の公式
\[
 \int_\alpha^\beta f(h(t))h'(t)\dd t
 =\int_{h(\alpha)}^{h(\beta)}f(x)\dd x
\]
自体も同じ合成形の公式であり,\,$h$ の単調性を必要としない.単調性が必要になるのは,逆関数を1つ選んで式を書き直す段階である.この2つを区別すれば,置換に単調性が要る場合と要らない場合を混同せずに済む.
